\section{Conclusão e trabalhos futuros}

Na base Adult, com \svLinhas{} pessoas e um teto de \svTeto\% imposto por perfis idênticos com rendas diferentes, a SVM com kernel RBF tem a melhor acurácia, \svRBFAcc~$\pm$~\svRBFSd\%: supera a SVM linear e a Random Forest nas 10 seeds, mas empata com a árvore ajustada, e a vantagem sobre uma fronteira reta é de menos de um ponto. O estudo numa divisão mostrou o que a SVM exige: padronização (sem ela, \svPadSem\%), um par de $C$ e $\gamma$ dentro de um platô largo, e um tempo de treino que cresce com o quadrado do tamanho do treino. A busca guiada pela acurácia escolhe um $\gamma$ alto, que dá à SVM RBF a pior AUC entre as SVMs (\svRBFAuc\%). A rede neural chega ao mesmo nível de acurácia (\svRedeAcc\%), com o melhor F1 macro e a melhor AUC (\svRedeAuc\%), uma busca cerca de cinco vezes mais curta e uma previsão centenas de vezes mais rápida. O melhor modelo depende, portanto, do critério: pela acurácia, a SVM RBF; pela capacidade de ordenar e pelo custo, a rede.

Os próximos passos seriam repetir a busca da SVM RBF com a AUC como critério, para medir quanto da diferença para a rede vem do critério e quanto do modelo; testar aproximações do kernel, que mantêm a fronteira curva com o custo de um modelo linear; e manter as linhas com valor ausente como uma categoria própria, para medir quanto a curadoria mudou os resultados.
